% META: {"id":"1SPE-EXPO-CO-009","chapitre":"1SPE-EXPONENTIELLE","type_objet":"corrige","exercice_id":"1SPE-EXPO-EX-009","status":"approved","extension_codes":["X4"],"programme_alignment":"OPTIONAL_EXTENSION","extension_label":"Approfondissement — Vers la Terminale","origin":"generated","status_history":["generated","audited","accepted","approved"],"release_acceptance":"RELEASE_OWNER_BATCH_ACCEPTANCE","acceptance_closure_digest":"sha256:9b3ccf9a81c5520fbb7e03b7b2d3e7bf2057a02834b6d908bf3be5f49f3b553f","acceptance_receipt_digest":"sha256:4fad86f0282e0fa4e0419cca1ad6379ae7af5a280c04a4a90880f90a1c044242"}
\begin{center}\textbf{Approfondissement — Vers la Terminale}\end{center}
% BEGIN-VERIFY
% from sympy import *
% x = symbols('x')
% f = Function('f')
% # g = f/exp, g' numerator = f'*exp - f*exp' = (f-f)*exp = 0 when f'=f and exp'=exp
% g_prime_num = diff(f(x), x) * exp(x) - f(x) * exp(x)
% g_prime_num_sub = g_prime_num.subs(diff(f(x), x), f(x))
% assert simplify(g_prime_num_sub) == 0
% # g(0) = f(0)/exp(0) = 1/1 = 1
% assert exp(0) == 1
% END-VERIFY
\begin{corrige}{1SPE-EXPO-EX-009}
\begin{enumerate}
\item La fonction $g(x) = \dfrac{f(x)}{\exp(x)}$ est définie à condition que $\exp(x)
\neq 0$. Or, pour tout réel $x$, $\exp(x) > 0$, donc $\exp(x)
\neq 0$. La fonction $g$ est bien définie sur $\mathbb{R}$.

\item Par la formule du quotient $\left(\dfrac{u}{v}\right)' = \dfrac{u'v - uv'}{v^2}$ avec $u = f$ et $v = \exp$ :
\[g'(x) = \frac{f'(x) \cdot \exp(x) - f(x) \cdot \exp'(x)}{[\exp(x)]^2}.\]

\item On utilise les hypothèses : $f'(x) = f(x)$ et $\exp'(x) = \exp(x)$. On remplace dans le numérateur :
\[f'(x) \cdot \exp(x) - f(x) \cdot \exp'(x) = f(x) \cdot \exp(x) - f(x) \cdot \exp(x) = 0.\]
Donc :
\[g'(x) = \frac{0}{[\exp(x)]^2} = 0 \quad \text{pour tout réel } x.\]

\item Une fonction dont la dérivée est nulle sur $\mathbb{R}$ est constante sur $\mathbb{R}$. Il existe donc une constante $C \in \mathbb{R}$ telle que $g(x) = C$ pour tout réel $x$.

\item On calcule $g(0)$ :
\[g(0) = \frac{f(0)}{\exp(0)} = \frac{1}{1} = 1.\]
Donc $C = 1$.

\item Pour tout réel $x$, $g(x) = 1$, c'est-à-dire :
\[\frac{f(x)}{\exp(x)} = 1, \quad \text{soit} \quad f(x) = \exp(x).\]

Toute fonction $f$ vérifiant $f' = f$ et $f(0) = 1$ est nécessairement égale à $\exp$. La fonction exponentielle est donc \textbf{l'unique} fonction dérivable sur $\mathbb{R}$ telle que $f' = f$ et $f(0) = 1$.
\end{enumerate}
\end{corrige}
